\documentclass[12pt]{article}

\usepackage{fontawesome}
\usepackage{hyperref}
\usepackage{xurl}
\usepackage[tagged, highstructure]{accessibility}
\usepackage{bookmark}
\usepackage{enumitem}

\hypersetup{
    colorlinks=false,
    pdfborder={0 0 0},
    pdftitle={Linux History},
    pdfauthor={Adrianna Holden-Gouveia},
    pdfsubject={Introduction to Linux - Lab Assignment},
    pdflang={en-US},
    pdfkeywords={lab assignment, Linux History}
}

% Configure PDF bookmarks for navigation
\bookmarksetup{
    numbered,
    open,
}

% Configure list spacing for better accessibility
\setlist{nosep}



\title{Linux History}
\author{
        Adrianna Holden-Gouveia \\
        Website: \href{https://aholdengouveia.name}{aholdengouveia.name}\\ 
        \date{\vspace{-5ex}}
        %Email: \href{mailto:admin@aholdengouveia.name}{admin@aholdengouveia.name} \\
%        \faLinkedin{: aholdengouveia} \\
        \faGithub {: aholdengouveia} \\
%        \faTwitter {: aholdengouveia} \\
        }

%S\date{\today}


\begin{document}    

\maketitle

\section*{Accessibility Notice}
This document is also available in HTML format at:

Link: \href{https://aholdengouveia.name/IntroLinux/labs/linuxhistory.html}{\textbf{https://aholdengouveia.name/IntroLinux/labs/linuxhistory.html}}

The HTML version provides enhanced accessibility features including keyboard navigation, screen reader support, responsive design, dark mode support, and high contrast options.


%\begin{abstract}
%This is a template for Linux Administration Lab work
%\end{abstract}
%\tableofcontents

\section*{Objectives:}
\begin{itemize}
    \item Learn about the history and significance of Linux in modern computing
    \item Develop effective note-taking skills using the sketchnotes format
    \item Practice summarizing technical content in your own words
\end{itemize}
\section*{Lab Instructions}

References, a video, a PowerPoint and some notes are available at my website
Link: \href{https://www.aholdengouveia.name/IntroLinux/linuxhistory.html}{\textbf{https://www.aholdengouveia.name/IntroLinux/linuxhistory.html}}


Read one of the following three articles and turn in a report using the format listed below.
It is very important that it is in your own words. Do not copy, plagiarize or slightly reword someone else's text.

You will be trying something called Sketchnotes.  There is more information here: Link: \href{https://rohdesign.com/sketchnotes}{\textbf{https://rohdesign.com/sketchnotes}}

For this lab you will read one of the articles and then make a page or less of notes, using the sketchnotes format, and turn in a picture of your notes.  Please make sure your picture is clear. Any paper and pen is fine, but please make sure it's legible.

\subsection*{Article options:}
    \begin{enumerate}
        \item Linux Took Over the Web. Now, It's Taking Over the World Link: \href{https://www.wired.com/2016/08/linux-took-web-now-taking-world/}{\textbf{https://www.wired.com/2016/08/linux-took-web-now-taking-world/}}
        \item Why Linux is the most important software project in history Link: \href{https://felipec.wordpress.com/2011/03/06/why-linux-is-the-most-important-software-project-in-history/}{\textbf{https://felipec.wordpress.com/2011/03/06/why-linux-is-the-most-important-software-project-in-history/}}
        \item Linux turns 33 in 2024 Link: \href{https://www.networkworld.com/article/1286507/linux-turns-33-in-2024.html}{\textbf{https://www.networkworld.com/article/1286507/linux-turns-33-in-2024.html}}

    \end{enumerate}



\section*{Deliverables}
Picture of your sketchnotes about the article you picked. Make sure somewhere on your notes you have your name and which article you picked. 
\end{document}